\chapter*{Danh mục các từ viết tắt}
\addcontentsline{toc}{chapter}{DANH MỤC CÁC TỪ VIẾT TẮT}
\begin{longtable}{@{}p{0.20\textwidth}p{0.75\textwidth}@{}}
\toprule
\textbf{Viết tắt} & \textbf{Ý nghĩa} \\
\midrule
\endfirsthead
\toprule
\textbf{Viết tắt} & \textbf{Ý nghĩa} \\
\midrule
\endhead
\midrule
\multicolumn{2}{r}{\textit{(Tiếp theo trang sau)}}\\
\endfoot
\bottomrule
\endlastfoot

AAD   & Additional Authenticated Data -- Dữ liệu xác thực bổ sung (trong cơ chế AEAD/AES-GCM) \\
ACK   & Acknowledgement -- Thông điệp xác nhận đã nhận/xử lý \\
AEAD  & Authenticated Encryption with Associated Data -- Mã hóa xác thực kèm dữ liệu liên kết \\
AES   & Advanced Encryption Standard -- Chuẩn mã hóa đối xứng AES \\
CORS  & Cross-Origin Resource Sharing -- Cơ chế chia sẻ tài nguyên khác nguồn (trình duyệt) \\
CSRF  & Cross-Site Request Forgery -- Tấn công giả mạo yêu cầu liên trang \\
CRUD  & Create, Read, Update, Delete -- 4 thao tác dữ liệu cơ bản \\
DB    & Database -- Cơ sở dữ liệu \\
E2EE  & End-to-End Encryption -- Mã hóa đầu-cuối \\
ESM   & ECMAScript Module -- Cơ chế module chuẩn của JavaScript \\
FIDO2 & Fast IDentity Online 2 -- Chuẩn xác thực không mật khẩu của FIDO Alliance \\
FIPS  & Federal Information Processing Standards -- Bộ tiêu chuẩn xử lý thông tin liên bang (Mỹ) \\
GCM   & Galois/Counter Mode -- Chế độ vận hành (mode) của AES cung cấp mã hóa xác thực \\
GDPR  & General Data Protection Regulation -- Quy định bảo vệ dữ liệu chung (EU) \\
GMAC  & Galois Message Authentication Code -- Mã xác thực dựa trên GCM (dạng MAC) \\
HMAC  & Hash-based Message Authentication Code -- Mã xác thực thông điệp dựa trên hàm băm \\
HMR   & Hot Module Replacement -- Cơ chế nạp nóng module trong môi trường phát triển \\
HS256 & HMAC-SHA-256 -- Thuật toán ký JWT dùng HMAC với SHA-256 \\
IV    & Initialization Vector -- Vector khởi tạo/nonce dùng cho mã hóa đối xứng (ngữ cảnh AES-GCM) \\
JSON  & JavaScript Object Notation -- Định dạng dữ liệu JSON \\
JWT   & JSON Web Token -- Token xác thực/phiên làm việc \\
MGF1  & Mask Generation Function 1 -- Hàm sinh mặt nạ (dùng trong OAEP/PSS) \\
NIST  & National Institute of Standards and Technology -- Viện tiêu chuẩn và công nghệ Hoa Kỳ \\
OAEP  & Optimal Asymmetric Encryption Padding -- Lược đồ đệm cho RSA dùng trong key wrapping \\
ODM   & Object--Document Mapper -- Ánh xạ Object--Document (MongoDB) \\
OTP   & One-Time Password -- Mật khẩu dùng một lần \\
PBKDF2& Password-Based Key Derivation Function 2 -- Hàm dẫn xuất khóa từ mật khẩu/PIN \\
PIN   & Personal Identification Number -- Mã PIN người dùng \\
PSS   & Probabilistic Signature Scheme -- Lược đồ chữ ký xác suất của RSA (RSA-PSS) \\
REST  & Representational State Transfer -- Phong cách thiết kế API (REST API) \\
RSA   & Rivest--Shamir--Adleman -- Thuật toán mật mã bất đối xứng RSA \\
RSAES & RSA Encryption Scheme -- Hệ lược đồ mã hóa RSA (ngữ cảnh RSAES-OAEP) \\
RSASSA& RSA Signature Scheme with Appendix -- Hệ lược đồ chữ ký RSA (ngữ cảnh RSASSA-PSS) \\
SHA   & Secure Hash Algorithm -- Họ hàm băm SHA (ngữ cảnh SHA-256) \\
SPA   & Single Page Application -- Ứng dụng web một trang \\
TLS   & Transport Layer Security -- Bảo mật tầng vận chuyển \\
TOFU  & Trust On First Use -- Tin cậy lần đầu; thay đổi khóa cần xác minh lại \\
UI    & User Interface -- Giao diện người dùng \\
UI/UX & User Interface / User Experience -- Giao diện người dùng / Trải nghiệm người dùng \\
URL   & Uniform Resource Locator -- Địa chỉ tài nguyên (đường dẫn) \\
USENIX& USENIX Association -- Tổ chức/hội nghị nghiên cứu hệ thống (ngữ cảnh trích dẫn) \\
W3C   & World Wide Web Consortium -- Tổ chức tiêu chuẩn web (W3C) \\
UX    & User Experience -- Trải nghiệm người dùng \\

\end{longtable}
